\documentclass[UTF8,zihao=-4,scheme=chinese]{ctexart}

\usepackage[
  a4paper,
  left=1.55cm,
  right=1.55cm,
  top=1.65cm,
  bottom=1.8cm,
  headheight=14pt,
  headsep=0.45cm,
  footskip=0.85cm
]{geometry}
\usepackage{setspace}
\setstretch{1.12}
\setlength{\parindent}{2em}
\setlength{\parskip}{0.25em}
\usepackage{amsmath,amssymb,mathtools}
\usepackage{enumitem}
\setlist[enumerate]{leftmargin=2.2em,itemsep=0.2em,topsep=0.25em}
\usepackage[hidelinks]{hyperref}
\usepackage{fancyhdr}
\pagestyle{fancy}
\fancyhf{}
\fancyhead[L]{\small 高等数学模拟卷}
\fancyhead[R]{\small 连续性专题}
\fancyfoot[C]{\small \thepage}
\renewcommand{\headrulewidth}{0.4pt}

\usepackage[most]{tcolorbox}
\tcbuselibrary{breakable,skins}
\definecolor{CAInk}{HTML}{1F2937}
\definecolor{CASubtle}{HTML}{F8FAFC}
\tcbset{
  ca-base/.style={
    enhanced,
    breakable,
    sharp corners,
    boxrule=0.6pt,
    leftrule=2.4pt,
    arc=1.2mm,
    left=1.1mm,
    right=1.1mm,
    top=0.9mm,
    bottom=0.9mm,
    before skip=0.65em,
    after skip=0.65em,
    fonttitle=\bfseries,
    coltitle=CAInk
  }
}
\newtcolorbox{solutionbox}[1][]{%
  ca-base,
  title={解答},
  colback=CASubtle,
  colframe=CAInk!55,
  colbacktitle=CAInk!8,
  #1
}

\title{\bfseries 高等数学连续性专题模拟卷}
\author{Course Alchemist 示例项目}
\date{时间：45 分钟 \qquad 满分：30 分}

\begin{document}

\maketitle

\section*{试题}

\begin{enumerate}
  \item（6 分）设
  \[
    f(x)=
    \begin{cases}
      x^2+c, & x<0,\\
      2x+1, & x\ge 0.
    \end{cases}
  \]
  求 $c$，使 $f$ 在 $x=0$ 处连续。

  \item（8 分）证明方程 $x^3+x-1=0$ 在区间 $(0,1)$ 内至少有一个实根。

  \item（8 分）设 $f$ 在 $[0,2]$ 上连续，且 $f(0)=1$，$f(2)=5$。证明存在 $\xi\in(0,2)$，使得 $f(\xi)=3$。

  \item（8 分）设 $f$ 在 $[0,1]$ 上连续，且 $0\le f(x)\le 1$。证明存在 $\xi\in[0,1]$，使得 $f(\xi)=\xi$。
\end{enumerate}

\section*{参考答案}

\begin{solutionbox}[title={第 1 题}]
连续要求
\[
  \lim_{x\to0^-}(x^2+c)=f(0).
\]
左极限为 $c$，而 $f(0)=1$，所以 $c=1$。
\end{solutionbox}

\begin{solutionbox}[title={第 2 题}]
令 $g(x)=x^3+x-1$。函数 $g$ 在 $[0,1]$ 上连续，且
\[
  g(0)=-1<0,\qquad g(1)=1>0.
\]
由零点存在定理，存在 $\xi\in(0,1)$，使得 $g(\xi)=0$。
\end{solutionbox}

\begin{solutionbox}[title={第 3 题}]
因为 $f$ 在 $[0,2]$ 上连续，且 $3$ 介于 $f(0)=1$ 与 $f(2)=5$ 之间，所以由介值定理，存在 $\xi\in(0,2)$，使得 $f(\xi)=3$。
\end{solutionbox}

\begin{solutionbox}[title={第 4 题}]
令 $g(x)=f(x)-x$。则 $g$ 在 $[0,1]$ 上连续，且
\[
  g(0)=f(0)\ge0,\qquad g(1)=f(1)-1\le0.
\]
若端点处 $g=0$，结论立即成立；否则由介值定理可得存在 $\xi\in(0,1)$，使得 $g(\xi)=0$，即 $f(\xi)=\xi$。
\end{solutionbox}

\end{document}
